\begin{resumo}
\SingleSpacing
\noindent
O modelo operacional de imobiliárias que gerenciam múltiplos corretores parceiros enfrenta um gargalo crítico no controle de acesso físico a propriedades dispersas. O compartilhamento tradicional de chaves metálicas e senhas de cofres mecânicos carece de rastreabilidade e apresenta riscos severos à segurança patrimonial. Embora fechaduras inteligentes conectadas mitiguem a distribuição de chaves, soluções vigentes pressupõem conectividade contínua à internet, inviabilizando sua aplicação no cenário mais frequente do setor: imóveis desocupados sem plano ativo de banda larga. Este trabalho propõe e avalia, sob a metodologia Design Science Research (DSR), um artefato de controle de acesso inteligente baseado em Internet das Coisas (IoT), validado em ambiente simulado, capaz de gerenciar com segurança o acesso temporário de múltiplos parceiros a imóveis desprovidos de conexão de rede. A solução integra controle de acesso baseado em atributos (ABAC contextual) a tokens criptográficos assinados digitalmente com o algoritmo de curva elíptica Ed25519, transmitidos via QR Code e verificados localmente pelo microcontrolador da fechadura em milissegundos, com sincronização assíncrona de auditoria via protocolo MQTT. A avaliação do artefato é estruturada por meio de simulação virtual em nós microcontrolados e testes automatizados de integração, estabelecendo métricas de latência temporal, ocupação de memória e eficácia de contenção lógica, com o intuito de fundamentar uma alternativa arquitetural de baixo custo e auditável no setor imobiliário.

\vspace{\onelineskip}
\noindent
\textbf{Palavras-chave}: Internet das Coisas (IoT). Controle de Acesso. Gestão Imobiliária. Criptografia Ed25519. Design Science Research.
\end{resumo}

\begin{resumo}[Abstract]
\SingleSpacing
\noindent
The operational model of real estate agencies managing multiple partner brokers faces a critical bottleneck in physical access control to dispersed properties. Traditional sharing of physical keys and lockbox passcodes lacks traceability and poses severe security risks. Although connected smart locks mitigate key handover logistics, current IoT solutions rely on continuous internet connectivity, failing in the sector's most common scenario: vacant properties without active broadband connections. This work proposes and evaluates, following the Design Science Research (DSR) methodology, an Internet of Things (IoT) intelligent access control artifact, validated in a simulated environment, capable of securely managing temporary access for multiple partners to offline properties. The solution couples Attribute-Based Access Control (contextual ABAC) with cryptographic tokens digitally signed via the Ed25519 elliptic curve algorithm, delivered through QR Codes and verified locally by the microcontrolled lock within milliseconds, alongside asynchronous audit synchronization via the MQTT protocol. The artifact evaluation is structured through virtual microcontroller simulation and automated integration tests, establishing performance metrics for latency, memory footprint, and logical containment efficacy, aiming to support an auditable and cost-effective architectural alternative tailored for real estate operations.

\vspace{\onelineskip}
\noindent
\textbf{Keywords}: Internet of Things (IoT). Access Control. Real Estate Management. Ed25519 Cryptography. Design Science Research.
\end{resumo}